\section{Методика выявления атак}

\subsection{Этап 1. Построение эталонного метаграфа угроз}

Эталонный метаграф \(G_{\text{ref}}\) строится на основе:
\begin{itemize}
    \item официальной базы MITRE ATT\&CK;
    \item реестра уязвимостей CVE/NVD;
    \item внутренних отчётов ИБ-команды;
    \item методики определения угроз \cite{fsteck2015threats};
    \item работы Новохрестова и Конева \cite{novokhrestov2014security}.
\end{itemize}

Каждая известная атака (например, APT28, Lazarus) представляется как подграф \(G_{\text{ref}}^{(i)}\), содержащий цепочку: цель → тактики → методы → уязвимости.

\subsection{Этап 2. Сбор и нормализация событий}

События поступают из разнородных источников:
\begin{itemize}
    \item EDR-агенты (запуск процессов, доступ к файлам);
    \item Сетевые IDS (аномальные потоки);
    \item Логи веб-серверов и баз данных.
\end{itemize}

Все события нормализуются в единый формат CEF или JSON, затем отображаются в метаграф наблюдений \(G_{\text{obs}}\).

\subsection{Этап 3. Алгоритм сопоставления}

Задача сводится к поиску \textbf{частичного изоморфизма} между \(G_{\text{ref}}^{(i)}\) и подграфом \(G_{\text{obs}}\). Используется рекурсивный алгоритм с эвристической фильтрацией по атрибутам \texttt{technique\_id} и \texttt{tactic\_id}.

\subsection{Этап 4. Оценка достоверности и реагирование}

Для гипотезы атаки рассчитывается вес:
\[
W(H_i) = \alpha \cdot S_{\text{struct}} + \beta \cdot S_{\text{attr}} + \gamma \cdot C_{\text{crit}},
\]
где \(S_{\text{struct}}\) — структурное совпадение, \(S_{\text{attr}}\) — совместимость атрибутов, \(C_{\text{crit}}\) — критичность цели. При \(W(H_i) > 0.85\) генерируется алерт.